% ============================================================================
% PHẦN 2 (2022 - ĐẦU 2023): CĂN CHỈNH MÔ HÌNH BẰNG HỌC TĂNG CƯỜNG (RLHF & PPO)
% ============================================================================
\sectionintro{2}{Alignment với RLHF và PPO}{2022--đầu 2023}{Sở thích cặp, Reward Model, KL regularization và kiến trúc bốn mạng.}

\subsection{2.1. Dữ liệu Đối chiếu Cặp \texorpdfstring{$(y_w, y_l)$}{(y\textunderscore w, y\textunderscore l)} \& Reward Model (\texorpdfstring{$r_\phi$}{r\textunderscore φ})}

% ----------------------------------------------------------------------------
% SLIDE: DỮ LIỆU ĐỐI CHIẾU CẶP: MẪU DƯƠNG y_w VÀ MẪU ÂM y_l
% ----------------------------------------------------------------------------
\begin{frame}{Học sở thích từ cặp phản hồi}
  \lead{InstructGPT chuyển từ mẫu đơn sang quan hệ thứ tự (Ouyang et al., 2022).}
  \subhead{Sinh phản hồi và đánh giá phản hồi}
  \small
  Thay vì soạn nhãn chuẩn $y^*$ như SFT, cho mô hình sinh cặp ứng viên $(y_1, y_2) \sim \pi_{\text{SFT}}(\cdot \mid x)$ rồi xác lập thứ tự $\mathbf{y_w \succ y_l}$ (\textbf{Mẫu dương $y_w$} ưu việt hơn \textbf{Mẫu âm $y_l$}):

  \vspace{0.08cm}
  \begin{columns}[T,onlytextwidth]
    \begin{column}{0.48\textwidth}
      \limithead{Vì sao không chỉ chấm điểm tuyệt đối?}
      \small
      \begin{itemize}\setlength{\itemsep}{4pt}
        \item \textbf{Phương sai cá nhân cao:} Thang điểm tuyệt đối ($1\text{--}10$) không đồng nhất giữa những người đánh giá.
        \item \textbf{Calibration Drift:} Ngưỡng chấm của cùng một chuyên gia biến thiên theo thời gian.
      \end{itemize}
    \end{column}
    \begin{column}{0.48\textwidth}
      \motivehead{Ưu thế của cặp đối chiếu $(y_w \succ y_l)$}
      \small
      \begin{itemize}\setlength{\itemsep}{4pt}
        \item \textbf{Bất biến với thang đo:} Quan hệ nhị phân $y_w \succ y_l$ có độ đồng thuận thống kê cao.
        \item \textbf{Bổ sung Mẫu âm $y_l$ (Negative Sample):} Cung cấp biên phân cách rõ ràng để triệt tiêu hành vi sai lệch.
      \end{itemize}
    \end{column}
  \end{columns}
  \par\vspace{0pt}
\end{frame}

% ----------------------------------------------------------------------------
% SLIDE: SƠ ĐỒ KIẾN TRÚC & LUỒNG DỮ LIỆU HUẤN LUYỆN REWARD MODEL
% ----------------------------------------------------------------------------
\begin{frame}{Huấn luyện Reward Model}
  \lead{Học điểm vô hướng $r_\phi(x,y)\in\mathbb{R}$ từ cặp mẫu dương và mẫu âm.}
  \vspace{0.15cm}
  \begin{center}
    \resizebox{0.96\textwidth}{!}{%
      \begin{tikzpicture}
        \node[data, text width=2.8cm, minimum height=1.2cm] (prompt) at (0, 0) {
          \textbf{Truy vấn $x$}\\
          (Từ tập $\mathcal{D}_{\text{pair}}$)
        };

        \node[data, fill=softgreen, draw=tealgreen, text width=3.3cm, minimum height=1.1cm] (yw) at (4.2, 1.15) {
          \textbf{Mẫu dương $y_w$}\\
          (Positive Sample)
        };

        \node[data, fill=softred, draw=stanfordred, text width=3.3cm, minimum height=1.1cm] (yl) at (4.2, -1.15) {
          \textbf{Mẫu âm $y_l$}\\
          (Negative Sample)
        };

        \node[trainmodel, text width=3.8cm, minimum height=3.1cm] (rm) at (9.0, 0) {
          \textbf{Reward Model $r_\phi$}\\
          (Reward Model)\\[6pt]
          \normalfont\scriptsize Khởi tạo từ $\pi_{\text{SFT}}$;\\
          Thay đầu từ vựng $|\mathcal{V}|$ bằng\\
          \textbf{Đầu Chiếu Vô hướng}\\
          $W_r \in \mathbb{R}^{d \times 1}$ tại \texttt{<EOS>}
        };

        \node[rewardbox, text width=2.9cm, minimum height=1.0cm] (rw) at (13.3, 1.15) {
          Điểm mẫu dương\\$\mathbf{r_\phi(x, y_w) \in \mathbb{R}}$
        };

        \node[lossbox, fill=softred, draw=stanfordred, text width=2.9cm, minimum height=1.0cm] (rl) at (13.3, -1.15) {
          Điểm mẫu âm\\$\mathbf{r_\phi(x, y_l) \in \mathbb{R}}$
        };

        \node[lossbox, text width=3.6cm, minimum height=1.6cm] (btloss) at (17.5, 0) {
          \textbf{Bradley--Terry Loss}\\[3pt]
          $\mathcal{L}_{\text{RM}} = -\log \sigma(r_w - r_l)$\\[3pt]
          \normalfont\scriptsize Cực đại hóa hiệu số $r_w - r_l$
        };

        \coordinate (bus_p) at (1.9, 0);
        \draw[line width=1.05pt, draw=navyblue!85!black] (prompt.east) -- (bus_p);
        \draw[flow] (bus_p) |- (yw.west);
        \draw[flow] (bus_p) |- (yl.west);

        \draw[flow] (yw.east) -- ([yshift=1.15cm]rm.west);
        \draw[flow] (yl.east) -- ([yshift=-1.15cm]rm.west);

        \draw[flow] ([yshift=1.15cm]rm.east) -- (rw.west);
        \draw[flow] ([yshift=-1.15cm]rm.east) -- (rl.west);

        \draw[flow] (rw.east) -| (btloss.north);
        \draw[flow] (rl.east) -| (btloss.south);

        \draw[gradflow] (btloss.south) -- (17.5, -2.35) -- node[below, font=\scriptsize\rmfamily\bfseries, text=stanfordred] {Backpropagation $\nabla_\phi \mathcal{L}_{\text{RM}}$ cập nhật tham số $\phi$} (9.0, -2.35) -- (rm.south);
      \end{tikzpicture}%
    }
  \end{center}
  \vspace{-0.05cm}
  \small
  \textbf{Nguyên lý vận hành:} Cặp $(x, y_w)$ và $(x, y_l)$ đi qua cùng mạng $r_\phi$; gradient đẩy điểm của mẫu dương $r_\phi(x, y_w)$ tăng lên ($\uparrow$) và điểm của mẫu âm $r_\phi(x, y_l)$ giảm xuống ($\downarrow$).
  \par\vspace{0pt}
\end{frame}

% ----------------------------------------------------------------------------
% SLIDE: TOÁN HỌC CỦA REWARD MODELING — MÔ HÌNH BRADLEY-TERRY & BẤT BIẾN TỊNH TIẾN
% ----------------------------------------------------------------------------
\begin{frame}{Mô hình Bradley--Terry và Loss}
  \lead{Maximum Likelihood Estimation cho quan hệ thứ tự cặp.}
  \begin{block}{Bradley--Terry (1952) và Loss $\mathcal{L}_{\text{RM}}(\phi)$}
    \small
    \[
      \fitmath{%
      P(y_w \succ y_l \mid x) = \sigma\big(r_\phi(x, y_w) - r_\phi(x, y_l)\big), \qquad
      \mathcal{L}_{\text{RM}}(\phi) = -\E_{(x, y_w, y_l) \sim \mathcal{D}_{\text{pair}}} \Big[ \log \sigma\big(r_\phi(x, y_w) - r_\phi(x, y_l)\big) \Big]
      }
    \]
  \end{block}

  \subhead{Bất biến tịnh tiến (Shift invariance)}
  \small
  \begin{itemize}\setlength{\itemsep}{4pt}
    \item \textbf{Biểu diễn vô hướng tại \texttt{<EOS>}:} $r_\phi(x, y) = W_r^\top h_{\text{EOS}}(x, y) \in \mathbb{R}$ (do mặt nạ Causal Attention hội tụ toàn bộ tiền tố tại token cuối).
    \item \textbf{Bất biến tịnh tiến:} Với mọi hàm $C(x)$ chỉ phụ thuộc truy vấn $x$, phép biến đổi $\tilde{r}(x, y) = r(x, y) + C(x)$ bảo toàn hiệu số:
      \[
        \fitmath{%
        \tilde{r}(x, y_w) - \tilde{r}(x, y_l) = r(x, y_w) - r(x, y_l) \implies P_{\tilde{r}}(y_w \succ y_l \mid x) \equiv P_r(y_w \succ y_l \mid x)
        }
      \]
  \end{itemize}
  \par\vspace{0pt}
\end{frame}

\subsection{2.2. Định luật Goodhart, KL regularization (\texorpdfstring{$\beta \DKL$}{β D\textunderscore KL}) \& Chu trình RL}

% ----------------------------------------------------------------------------
% SLIDE: ĐỊNH LUẬT GOODHART & REWARD HACKING -> ĐỘNG CƠ PHẠT KL
% ----------------------------------------------------------------------------
\begin{frame}{Vì sao cần phạt KL divergence?}
  \lead{Tối đa hóa trực tiếp $\E[r_\phi(x,y)]$ có thể làm suy biến phân phối.}
  \begin{columns}[T,onlytextwidth]
    \begin{column}{0.48\textwidth}
      \limithead{Goodhart và Reward Hacking}
      \small
      \begin{itemize}\setlength{\itemsep}{4pt}
        \item \textbf{Định luật Goodhart:} Khi hàm xấp xỉ $r_\phi$ bị tối ưu hóa cực đoan, nó mất tính tương quan với chất lượng thực tế.
        \item \textbf{OOD Error:} $r_\phi$ chỉ chính xác trong lân cận phân phối của $\pi_{\text{SFT}}$.
        \item \textbf{Reward Hacking (Gao et al., 2023):} Policy $\pi_\theta$ dịch chuyển sang vùng ngoại lai để khai thác cực đại giả của $r_\phi$.
      \end{itemize}
    \end{column}
    \begin{column}{0.48\textwidth}
      \motivehead{Điều chuẩn vùng tin cậy bằng $\DKL$}
      \small
      \begin{itemize}\setlength{\itemsep}{4pt}
        \item Cố định mô hình SFT làm \textbf{Reference Model ($\pi_{\text{ref}}$)} và phạt độ lệch phân phối:
          \[
            \fitmath{%
            \DKL\big(\pi_\theta(\cdot \mid x) \parallel \pi_{\text{ref}}(\cdot \mid x)\big) = \E_{y \sim \pi_\theta} \left[ \log \frac{\pi_\theta(y \mid x)}{\pi_{\text{ref}}(y \mid x)} \right]
            }
          \]
        \item \textbf{Tác dụng:} Neo $\pi_\theta$ trong vùng tin cậy của $r_\phi$ và bảo toàn cấu trúc ngôn ngữ gốc.
      \end{itemize}
    \end{column}
  \end{columns}
  \par\vspace{0pt}
\end{frame}

% ----------------------------------------------------------------------------
% SLIDE: CÔNG THỨC TOÁN HỌC CỦA RLHF CÓ RÀNG BUỘC PHẠT KL
% ----------------------------------------------------------------------------
\begin{frame}{RLHF Objective}
  \lead{Tối đa hóa phần thưởng và KL regularization (Ouyang et al., 2022).}
  \begin{block}{Mục tiêu RLHF có KL regularization}
    \[
      \fitmath{%
      \max_{\pi_\theta} \; J_{\text{RLHF}}(\pi_\theta) = \E_{x \sim \mathcal{D},\; y \sim \pi_\theta(y \mid x)} \Big[ \underbrace{r_\phi(x, y)}_{\text{Điểm định giá } r_\phi} - \;\beta \underbrace{\log \frac{\pi_\theta(y \mid x)}{\pi_{\text{ref}}(y \mid x)}}_{\text{KL regularization}} \Big]
      }
    \]
  \end{block}

  \subhead{Phần thưởng hiệu dụng và hệ số $\beta > 0$}
  \small
  \begin{itemize}\setlength{\itemsep}{4pt}
    \item Phần thưởng tổng hợp tại mỗi bước lấy mẫu trực tuyến $y \sim \pi_\theta(\cdot \mid x)$:
      \[
        \fitmath{%
        R_{\text{total}}(x, y) = r_\phi(x, y) - \beta \log \frac{\pi_\theta(y \mid x)}{\pi_{\text{ref}}(y \mid x)}
        }
      \]
    \item \textbf{Giới hạn tiệm cận:} Khi $\beta \to 0 \implies$ suy biến do \textit{Reward Hacking}; khi $\beta \to +\infty \implies$ Policy bất động ($\pi_\theta \equiv \pi_{\text{ref}}$).
  \end{itemize}
  \par\vspace{0pt}
\end{frame}

% ----------------------------------------------------------------------------
% SLIDE: SƠ ĐỒ Online RL Loop & KHẮC PHỤC EXPOSURE BIAS
% ----------------------------------------------------------------------------
\begin{frame}{Online RL Loop}
  \lead{Policy tự sinh mẫu $\hat y\sim\pi_\theta$ rồi nhận phản hồi để cập nhật.}
  \vspace{0.12cm}
  \begin{center}
    \resizebox{0.96\textwidth}{!}{%
      \begin{tikzpicture}
        \node[data, text width=3.1cm, minimum height=1.3cm] (state) at (0, 0) {
          \textbf{Trạng thái $s_t$}\\
          Truy vấn $x$ \& tiền tố\\tự sinh $\hat{y}_{<t}$
        };

        \node[trainmodel, text width=3.5cm, minimum height=1.3cm] (policy) at (4.6, 0) {
          \textbf{Policy $\pi_\theta(\hat{y}_t \mid s_t)$}\\
          Mô hình Ngôn ngữ\\đang tối ưu hóa
        };

        \node[data, fill=softgreen, draw=tealgreen, text width=3.4cm, minimum height=1.3cm] (action) at (9.3, 0) {
          \textbf{Quỹ đạo tự sinh}\\
          Rút mẫu $\hat{y} \sim \pi_\theta(\cdot \mid x)$\\
          (On-Policy Rollout)
        };

        \node[rewardbox, text width=3.5cm, minimum height=1.3cm] (env) at (14.0, 0) {
          \textbf{Định giá Phần thưởng}\\
          $R_{\text{total}} = r_\phi(x, \hat{y}) - \beta \DKL$
        };

        \node[lossbox, text width=7.8cm, minimum height=1.1cm] (update) at (9.3, -2.0) {
          \textbf{Advantage Function $A(x, \hat{y}) = R_{\text{total}}(x, \hat{y}) - b(x)$}\\[2pt]
          \normalfont\scriptsize Khi $A > 0$: Tăng xác suất $\pi_\theta(\hat{y} \mid x)$ ($\uparrow$) \quad $|$ \quad Khi $A < 0$: Giảm xác suất $\pi_\theta(\hat{y} \mid x)$ ($\downarrow$)
        };

        \draw[flow] (state.east) -- (policy.west);
        \draw[flow] (policy.east) -- (action.west);
        \draw[flow] (action.east) -- (env.west);

        \draw[flow] (env.south) |- (update.east);
        \draw[gradflow] (update.west) -| node[pos=0.25, below, font=\scriptsize\rmfamily\bfseries, text=stanfordred] {Cập nhật Policy Gradient $\nabla_\theta J(\theta)$} (policy.south);
      \end{tikzpicture}%
    }
  \end{center}
  \vspace{0.05cm}
  \small
  \textbf{Khắc phục Exposure Bias:} Policy $\pi_\theta$ được huấn luyện trực tiếp trên các trạng thái tự sinh $\hat{s}_t = (x, \hat{y}_{<t}) \sim \pi_\theta$ và nhận phạt âm ($A < 0$) ngay khi đi chệch quỹ đạo.
  \par\vspace{0pt}
\end{frame}

\subsection{2.3. Policy Gradient, Thuật toán PPO (Value Network \texorpdfstring{$V_\psi$}{V\textunderscore ψ}) \& Hệ thống 4 Mô hình}

% ----------------------------------------------------------------------------
% SLIDE: ĐỊNH LÝ POLICY GRADIENT & baseline BASELINE
% ----------------------------------------------------------------------------
\begin{frame}{Policy Gradient và baseline}
  \lead{REINFORCE xử lý phép lấy mẫu token rời rạc (Williams, 1992).}
  \begin{block}{REINFORCE với baseline và Advantage $A(x, y)$}
    \[
      \fitmath{%
      \nabla_\theta J(\theta) = \E_{x \sim \mathcal{D},\; y \sim \pi_\theta(\cdot \mid x)} \Big[ \underbrace{\big(R_{\text{total}}(x, y) - b(x)\big)}_{\text{Advantage Function } A(x, y)} \, \nabla_\theta \log \pi_\theta(y \mid x) \Big]
      }
    \]
  \end{block}

  \begin{columns}[T,onlytextwidth]
    \begin{column}{0.49\textwidth}
      \subhead{Đồng nhất thức gradient logarit}
      \small
      \begin{itemize}\setlength{\itemsep}{4pt}
        \item Phép rút mẫu rời rạc $y \sim \pi_\theta$ không có gradient $\frac{\partial y}{\partial \theta}$.
        \item Áp dụng $\nabla_\theta \pi_\theta = \pi_\theta \nabla_\theta \log \pi_\theta$:
          \[
            \fitmath{%
            \nabla_\theta \sum_y \pi_\theta(y) R(y) = \E_{y \sim \pi_\theta}\big[R(y)\nabla_\theta \log \pi_\theta(y)\big]
            }
          \]
      \end{itemize}
    \end{column}
    \begin{column}{0.48\textwidth}
      \motivehead{Vai trò của baseline $b(x) = \E[R]$}
      \small
      \begin{itemize}\setlength{\itemsep}{4pt}
        \item Phân tách dấu giữa quỹ đạo trên trung bình ($A > 0, \uparrow$) và dưới trung bình ($A < 0, \downarrow$).
        \item Giảm phương sai ước lượng Monte Carlo mà không làm chệch kỳ vọng (do $\E[b(x)\nabla_\theta \log \pi_\theta] = 0$).
      \end{itemize}
    \end{column}
  \end{columns}
  \par\vspace{0pt}
\end{frame}

% ----------------------------------------------------------------------------
% SLIDE: THUẬT TOÁN PPO: Critic & CƠ CHẾ PPO-CLIP
% ----------------------------------------------------------------------------
\begin{frame}{PPO: Clipping và Value Estimation}
  \lead{Kết hợp PPO-Clip và Critic (Schulman et al., 2017).}
  \begin{block}{PPO-Clip Objective}
    \[
      \fitmath{%
      \mathcal{J}_{\text{PPO-Clip}}(\theta) = \E_{t} \Big[ \min \Big( \rho_t(\theta) \hat{A}_t, \; \clip\big(\rho_t(\theta), 1 - \epsilon, 1 + \epsilon\big) \hat{A}_t \Big) \Big],
      \qquad \rho_t(\theta) = \frac{\pi_\theta(y_t \mid x, y_{<t})}{\pi_{\theta_{\text{old}}}(y_t \mid x, y_{<t})}
      }
    \]
  \end{block}

  \begin{columns}[T,onlytextwidth]
    \begin{column}{0.49\textwidth}
      \subhead{Clipping $\clip(\rho_t, 1-\epsilon, 1+\epsilon)$}
      \small
      \begin{itemize}\setlength{\itemsep}{4pt}
        \item $\rho_t(\theta)$: Tỷ số xác suất của token $y_t$ giữa Policy mới và Policy lúc lấy mẫu.
        \item Khi $|\rho_t(\theta) - 1| > \epsilon$ (với $\epsilon = 0.2$), toán tử $\min(\cdot)$ triệt tiêu gradient về $0$, ngăn bước cập nhật quá lớn.
      \end{itemize}
    \end{column}
    \begin{column}{0.47\textwidth}
      \subhead{Critic $V_\psi(s_t)$}
      \small
      \begin{itemize}\setlength{\itemsep}{4pt}
        \item Ước lượng baseline kỳ vọng tại từng vị trí token $t$: $\hat{A}_t = R_t - V_\psi(s_t)$.
        \item Đòi hỏi duy trì một mạng nơ-ron $V_\psi$ có quy mô tham số tương đương mô hình chính $\pi_\theta$.
      \end{itemize}
    \end{column}
  \end{columns}
  \par\vspace{0pt}
\end{frame}

% ----------------------------------------------------------------------------
% SLIDE: SƠ ĐỒ KIẾN TRÚC 4 MÔ HÌNH CỦA RLHF (PPO)
% ----------------------------------------------------------------------------
\begin{frame}{RLHF + PPO: bốn mô hình phối hợp}
  \lead{Actor, Reference, Reward Model và Critic cùng hoạt động trên GPU.}
  \vspace{0.08cm}
  \begin{center}
    \resizebox{0.96\textwidth}{!}{%
      \begin{tikzpicture}
        \node[data, text width=2.4cm, minimum height=1.2cm] (prompt) at (0, 0) {
          \textbf{Truy vấn $x$}\\
          (Từ tập dữ liệu)
        };

        \node[trainmodel, text width=3.5cm, minimum height=1.35cm] (actor) at (4.2, 1.35) {
          \textbf{1. Actor / Policy $\pi_\theta$}\\
          (Đang huấn luyện)\\
          \normalfont\scriptsize Tự sinh chuỗi $y \sim \pi_\theta(\cdot \mid x)$
        };

        \node[frozenmodel, text width=3.5cm, minimum height=1.35cm] (ref) at (4.2, -1.35) {
          \textbf{2. Reference $\pi_{\text{ref}}$}\\
          (Cố định tham số)\\
          \normalfont\scriptsize Tính $\log \pi_{\text{ref}}(y \mid x)$
        };

        \node[data, fill=softgreen, draw=tealgreen, text width=2.5cm, minimum height=1.2cm] (resp) at (8.5, 0) {
          \textbf{Chuỗi tự sinh $y$}\\
          (On-policy Rollout)
        };

        \node[frozenmodel, text width=3.6cm, minimum height=1.35cm] (rm) at (12.8, 1.35) {
          \textbf{3. Reward Model $r_\phi$}\\
          (Cố định tham số)\\
          \normalfont\scriptsize Định giá $r_\phi(x, y)$
        };

        \node[trainmodel, fill=softyellow, draw=amberorange!90!black, text width=3.6cm, minimum height=1.35cm] (critic) at (12.8, -1.35) {
          \textbf{4. Critic / Value $V_\psi$}\\
          (Đang huấn luyện)\\
          \normalfont\scriptsize Ước lượng nền $V_\psi(s_t)$
        };

        \node[lossbox, text width=3.7cm, minimum height=2.2cm] (ppo) at (17.5, 0) {
          \textbf{Tối ưu hóa PPO}\\[3pt]
          $R_{\text{total}} = r_\phi - \beta \DKL$\\[2pt]
          $\hat{A}_t = R_{\text{total}} - V_\psi(s_t)$\\[3pt]
          \normalfont\scriptsize Cắt tỉa tỷ số $\clip(\rho_t)$
        };

        \coordinate (bus_p) at (1.8, 0);
        \draw[line width=1.05pt, draw=navyblue!85!black] (prompt.east) -- (bus_p);
        \draw[flow] (bus_p) |- (actor.west);
        \draw[flow] (bus_p) |- (ref.west);

        \draw[flow] (actor.east) -| (resp.north);
        \draw[flow] (resp.south) |- (ref.east);

        \coordinate (bus_r) at (10.3, 0);
        \draw[line width=1.05pt, draw=navyblue!85!black] (resp.east) -- (bus_r);
        \draw[flow] (bus_r) |- (rm.west);
        \draw[flow] (bus_r) |- (critic.west);

        \draw[flow] (rm.east) -| ([xshift=-0.4cm]ppo.north);
        \draw[flow] (critic.east) -| ([xshift=-0.4cm]ppo.south);
        \draw[flow] (ref.south) -- (4.2, -2.45) -- node[below, font=\scriptsize\rmfamily, text=slate] {Truyền $\log \pi_{\text{ref}}(y \mid x)$ để tính phân kỳ $\beta \DKL(\pi_\theta \parallel \pi_{\text{ref}})$} (17.5, -2.45) -- (ppo.south);

        \draw[gradflow] (ppo.north) -- (17.5, 2.45) -- node[above, font=\scriptsize\rmfamily\bfseries, text=stanfordred] {Cập nhật Gradient PPO đồng thời cho Policy $\pi_\theta$ và Critic $V_\psi$} (4.2, 2.45) -- (actor.north);
      \end{tikzpicture}%
    }
  \end{center}
  \par\vspace{0pt}
\end{frame}

\subsection{2.4. Tổng kết Phần 2 (RLHF + PPO): Thành tựu \& Các Vấn đề Tồn đọng cho Phần 3}

% ----------------------------------------------------------------------------
% SLIDE: TỔNG KẾT PHẦN 2 (1/2): NHỮNG GÌ ĐÃ ĐẠT ĐƯỢC & TỐI ƯU SO VỚI SFT
% ----------------------------------------------------------------------------
\begin{frame}{RLHF + PPO cải thiện SFT như thế nào?}
  \lead{Đối chiếu RLHF + PPO (2022) và SFT (2021).}
  \motivehead{Đóng góp so với SFT}
  \vspace{0.05cm}
  \footnotesize
  \renewcommand{\arraystretch}{1.20}
  \begin{tabularx}{\textwidth}{@{}p{3.4cm} X X@{}}
    \toprule
    \textbf{Tiêu chí đánh giá} & \textbf{SFT (Phần 1 -- 2021)} & \textbf{RLHF + PPO (Phần 2 -- 2022)} \\
    \midrule
    \textbf{1. Phân phối trạng thái} & Học tĩnh trên tiền tố chuyên gia $y_{<t}^*$ (\textit{Teacher Forcing}) & \textbf{\textcolor{tealgreen}{Khử Exposure Bias}}: Học on-policy trên chuỗi tự sinh $\hat{y} \sim \pi_\theta$ (\textit{On-policy}) \\

    \textbf{2. Tín hiệu gradient} & Chỉ có tín hiệu dương ($\uparrow$) trên mẫu chuẩn $y^*$ & \textbf{\textcolor{tealgreen}{Gradient hai chiều}}: Củng cố mẫu dương ($\hat{A}_t > 0, \uparrow$) và chế tài mẫu âm ($\hat{A}_t < 0, \downarrow$) \\

    \textbf{3. Kiểm soát phân phối} & Suy giảm hiệu chỉnh xác suất khi ép học thuộc & Điều chuẩn vùng tin cậy $\beta\DKL(\pi_\theta \parallel \pi_{\text{ref}})$ chống suy biến \\

    \textbf{4. Hiệu năng thực nghiệm} & Phụ thuộc mạnh vào kích thước Base Model & \textbf{InstructGPT 1.3B} vượt qua \textbf{GPT-3 175B} trên đánh giá mù \\
    \bottomrule
  \end{tabularx}
  \par\vspace{0pt}
\end{frame}

% ----------------------------------------------------------------------------
% SLIDE: TỔNG KẾT PHẦN 2 (2/2): VẤN ĐỀ TỒN ĐỌNG MÀ PHẦN 3 (DPO) GIẢI QUYẾT
% ----------------------------------------------------------------------------
\begin{frame}{Các điểm nghẽn của RLHF + PPO}
  \lead{Chi phí bộ nhớ, thời gian sinh mẫu và tính ổn định dẫn đến DPO (2023).}
  \limithead{Ba điểm nghẽn hệ thống}
  \small
  \begin{enumerate}\setlength{\itemsep}{4pt}
    \item \textbf{Chi phí Bộ nhớ GPU gấp 4 lần (4-Model Memory Bottleneck):}\\
      \small Cần nạp đồng thời \textbf{4 mạng nơ-ron lớn} ($\pi_\theta, \pi_{\text{ref}}, r_\phi, V_\psi$) vào VRAM.
    \item \textbf{Thời gian Autoregressive Decoding (Autoregressive Rollout Bottleneck):}\\
      \small Sinh mẫu tuần tự $\hat{y} \sim \pi_\theta$ tại mỗi bước cập nhật tiêu tốn $>70\%$ thời gian tính toán.
    \item \textbf{Phương sai gradient cao và sai số xấp xỉ kép ($r_\phi, V_\psi$):}\\
      \small Sai số ước lượng từ cả Reward Model $r_\phi$ lẫn Value Network $V_\psi$ gây mất ổn định hội tụ.
  \end{enumerate}

  \vspace{0.06cm}
  \begin{exampleblock}{Tiếp theo: DPO (2023)}
    \small
    Dùng \textbf{closed-form solution} của RLHF để cập nhật trực tiếp $\pi_\theta$ từ cặp mẫu dương/âm $(y_w, y_l)$ mà \textbf{bỏ hoàn toàn vòng lặp RL, bỏ cả $r_\phi$ lẫn $V_\psi$}.
  \end{exampleblock}
  \par\vspace{0pt}
\end{frame}
